\scanpage{5}
4) \uline{Lemme} Soient $M,N$ des groupes commutatifs, et $\Sigma\colon M\longrightarrow {}_2N$ un hom. Alors il existe une application bilinéaire antisymétrique $\varphi\colon M\times M\longrightarrow N$ telle que $\varphi(x,x)=\Sigma(x)$ pour tt $x\in M$.

Remplaçant $M,N$ resp. par $\deleted{\unsure{\Sigma|2}}M_2,{}_2N$, on peut supposer $M,N$ des vectoriels \deleted{\unsure{de dim}} sur $\mathbb F_2$.

Soit $(e_i)_{i\in I}$ une base de $M$, \deleted{\unsure{on}} Alors $\deleted{\varphi}\inserted{\Sigma}$ est donnée par les $n_i=\Sigma(e_i)\in N$, tandis que $\varphi$ sera donnée par les $\varphi(e_i,e_j)=n_{ij}\in N$; la condition d'antisymétrie \inserted{(= symétrie)} est que $\deleted{\unsure{\varphi(e_i,e_j)=}} n_{ij}=n_{ji}$. On prendra $\varphi(e_i,\unsure{e_j|j})=\begin{cases}0&\text{si }i\ne j\\ n_i&\text{si }i=j\end{cases}$ \unsure{O.K.|OK}

\inserted{\uline{Corollaire} Toute catégorie de Picard devient \uline{stricte} en changeant sa seule contrainte de commutativité. En tant que Gr-catégorie, elle est déjà triviale\ldots}

5) Soit $C$ une catégorie de Picard, $\Gamma$ un groupe commutatif, \deleted{$\varphi$}
\[
\varphi\colon\Gamma\times\Gamma\longrightarrow\pi_1(C)
\]
une application bilinéaire antisymétrique. On désigne par $\Gamma\overset{\varphi}{\times}C$ la catégorie de Picard déduite de $\Gamma\times C$ (cf n° 2) par l'application bilinéaire \unsure{antisym|antisym.}
\[
\bar\varphi\colon\pi_0(\Gamma\times C)\times\pi_0(\Gamma\times C)\longrightarrow\pi_1(\Gamma\times C)
\]
(cf n° 3) donnée par
\[
\bar\varphi\big((\alpha,x),(\beta,y)\big)=\varphi(\alpha,\beta)
\]
i.e. composée \deleted{\unsure{de}}
\diagram{\begin{tikzcd}[ampersand replacement=\&,column sep=small]
{\pi_0(\Gamma\times C)\times\pi_0(\Gamma\times C)} \arrow[dr] \& \& {\pi_1(\Gamma\times C)\simeq\pi_1(C)}\\
\& {\deleted{\unsure{\pi_0(}}\Gamma\times\Gamma} \arrow[ur,bend left=10,draw=none,"\deleted{\longrightarrow}"{description,sloped}] \arrow[ur,bend right=7,"\varphi"'] \&
\end{tikzcd}}

6)
% archive sheet number: 4
